\section{Análise de Resultados}
O pipeline processou com sucesso 50 arquivos brutos, resultando em uma base consolidada que permite consultas instantâneas.

\subsection{Métricas de Processamento}
A automação reduziu o tempo de coleta de dados de uma estimativa manual de 2 horas para menos de 10 minutos de processamento computacional ininterrupto. O banco de dados populado permite agora a filtragem de qualquer um dos municípios brasileiros em tempo real.

\subsection{Visualização Analítica (Dashboard)}
O dashboard desenvolvido em Streamlit cumpre os requisitos de exibição de dados brutos e estatísticas descritivas. Através de bibliotecas como Plotly, foram gerados gráficos de tendência que evidenciam a sazonalidade da produção ambulatorial ao longo dos 25 meses analisados.

\subsection{Análise Descritiva}
Observou-se que a variável "Quantidade Aprovada" apresenta uma distribuição altamente assimétrica, concentrada em grandes centros urbanos, o que valida a necessidade de filtros por região para uma análise de saúde pública mais equânime.
